%======================================================================
% پیوست الف: مستندات API
%======================================================================
\chapter{مستندات \lr{API}}
\label{app:api}

در این پیوست، قرارداد بیرونی سامانه به‌صورت منظم مستند شده است. فهرست مسیرها
مستقیماً از فایل مسیرها، کنترلرها، منبع‌های پاسخ و آزمون‌های خودکار پروژه
استخراج شده است؛ بنابراین این پیوست نه یک توصیف آرمانی، بلکه شرحی از
رفتار واقعی نسخه‌ی پیاده‌سازی‌شده است.

\section{قراردادهای مشترک}

\subsection*{احراز هویت و سربرگ‌ها}
مسیرهای کاربری و مدیریتی مانند \code{/me}، \code{/wallet}،
\code{/plans} و \code{/reports} از توکن \lr{Sanctum} کاربر استفاده
می‌کنند و قالب رایج سربرگ آن‌ها
\code{Authorization: Bearer <SANCTUM\_TOKEN>} است. مسیرهای هوش مصنوعی از
کلید پروژه استفاده می‌کنند و دو قالب
\code{Authorization: Bearer <API\_KEY>} و \code{X-API-Key: <API\_KEY>} را
می‌پذیرند. همچنین مسیر \code{GET /api/v1/invoices/\{invoice\}/pdf} بر
پایه‌ی پیوند امضاشده عمل می‌کند و معمولاً از فیلد \code{pdf\_url} در پاسخ
فاکتور به‌دست می‌آید.

\subsection*{نحوه‌ی بازگشت \lr{Resource Collection}}
در این پروژه \code{JsonResource::withoutWrapping()} فعال شده است؛ بنابراین
پاسخ‌های مجموعه‌ای با لفافه‌ی \code{data} برگردانده نمی‌شوند و مستقیماً به
صورت آرایه‌ی \lr{JSON} در خروجی قرار می‌گیرند.

\subsection*{الگوی خطا در مسیرهای عمومی}
مسیرهای عمومی و مدیریتی خطاها را معمولاً با یک فیلد ساده‌ی \code{message}
بازمی‌گردانند.

\begin{latin}
\begin{lstlisting}[style=jsonStyle]
{
  "message": "The name field is required."
}
\end{lstlisting}
\end{latin}

\listingcaption{نمونه خطا در مسیرهای عمومی}

\subsection*{الگوی خطا در مسیرهای هوش مصنوعی}
درگاه هوش مصنوعی برای سازگاری بیرونی، خطاها را در قالبی مشابه \lr{OpenAI}
بازمی‌گرداند.

\begin{latin}
\begin{lstlisting}[style=jsonStyle]
{
  "error": {
    "message": "Provider and model are required.",
    "type": "invalid_request_error",
    "param": "provider",
    "code": null
  }
}
\end{lstlisting}
\end{latin}

\listingcaption{نمونه خطا در مسیرهای هوش مصنوعي}

\section{نقاط پایانی احراز هویت}

\subsection{\lr{POST /api/v1/auth/otp/start}}
\textbf{کارکرد.} آغاز فرایند احراز هویت با ارسال رمز یک‌بارمصرف.

\textbf{احراز هویت.} نیاز ندارد.

\textbf{ساختار درخواست.} بدنه‌ی درخواست شامل دو فیلد اجباری
\code{channel} و \code{destination} است. مقدار \code{channel} فقط می‌تواند
\code{email} یا \code{sms} باشد.

\textbf{ساختار پاسخ موفق.} یک شیء ساده با فیلد \code{status} که در حالت
موفق مقدار آن \code{sent} است.

\begin{latin}
\begin{lstlisting}[style=jsonStyle]
POST /api/v1/auth/otp/start
Content-Type: application/json

{
  "channel": "email",
  "destination": "user@example.com"
}
\end{lstlisting}
\end{latin}

\listingcaption{نمونه درخواست \code{OTP Start}}

\begin{latin}
\begin{lstlisting}[style=jsonStyle]
{
  "status": "sent"
}
\end{lstlisting}
\end{latin}

\listingcaption{نمونه پاسخ \code{OTP Start}}

\subsection{\lr{POST /api/v1/auth/otp/verify}}
\textbf{کارکرد.} تأیید رمز یک‌بارمصرف و صدور توکن دسترسی داشبورد.

\textbf{احراز هویت.} نیاز ندارد.

\textbf{ساختار درخواست.} فیلدهای اجباری \code{destination} و \code{code}
و فیلد اختیاری \code{channel} دریافت می‌شوند.

\textbf{ساختار پاسخ موفق.} شیئی شامل \code{token}، \code{token\_type}
و \code{user}. شیء \code{user} دارای فیلدهای \code{id}، \code{name}،
\code{email}، \code{phone}، \code{status}، \code{profile\_image\_url}،
\code{created\_at} و \code{updated\_at} است.

\begin{latin}
\begin{lstlisting}[style=jsonStyle]
POST /api/v1/auth/otp/verify
Content-Type: application/json

{
  "destination": "user@example.com",
  "code": "123456",
  "channel": "email"
}
\end{lstlisting}
\end{latin}

\listingcaption{نمونه درخواست \code{OTP Verify}}

\begin{latin}
\begin{lstlisting}[style=jsonStyle]
{
  "token": "1|wQy0fP4...sanctum-token",
  "token_type": "Bearer",
  "user": {
    "id": 10,
    "name": "User Example",
    "email": "user@example.com",
    "phone": null,
    "status": "active",
    "profile_image_url": null,
    "created_at": "2026-01-29T08:00:00.000000Z",
    "updated_at": "2026-01-29T08:00:00.000000Z"
  }
}
\end{lstlisting}
\end{latin}

\listingcaption{نمونه پاسخ \code{OTP Verify}}

\section{نقاط پایانی کاربر}

\subsection{\lr{GET /api/v1/me}}
\textbf{کارکرد.} دریافت پروفایل کاربر جاری.

\textbf{احراز هویت.} توکن \lr{Sanctum}.

\textbf{ساختار درخواست.} این مسیر بدنه ندارد.

\textbf{ساختار پاسخ موفق.} پاسخ این مسیر یک شیء کاربر است.
فیلدهای اصلی آن \code{id}، \code{name}، \code{email}، \code{phone}،
\code{status}، \code{profile\_image\_url}، \code{created\_at} و
\code{updated\_at} هستند.

\subsection{\lr{PATCH /api/v1/me}}
\textbf{کارکرد.} به‌روزرسانی پروفایل کاربر.

\textbf{احراز هویت.} توکن \lr{Sanctum}.

\textbf{ساختار درخواست.} فیلدهای \code{name}، \code{email} و \code{phone}
اختیاری هستند. در صورت ارسال تصویر، فیلد \code{profile\_image} باید به‌صورت
فایل تصویری و در قالب \lr{multipart/form-data} ارسال شود.

\textbf{ساختار پاسخ موفق.} همان ساختار مسیر \code{GET /api/v1/me} با
مقادیر به‌روزشده.

\section{نقاط پایانی کلیدهای \lr{API}}

\subsection{\lr{GET /api/v1/api-clients}}
\textbf{کارکرد.} دریافت فهرست مشتریان \lr{API} متعلق به کاربر جاری.

\textbf{احراز هویت.} توکن \lr{Sanctum}.

\textbf{ساختار درخواست.} این مسیر بدنه ندارد.

\textbf{ساختار پاسخ موفق.} آرایه‌ای از اشیا با فیلدهای \code{id}،
\code{name}، \code{status}، \code{created\_at} و \code{updated\_at}.

\subsection{\lr{POST /api/v1/api-clients}}
\textbf{کارکرد.} ایجاد مشتری \lr{API} جدید.

\textbf{احراز هویت.} توکن \lr{Sanctum}.

\textbf{ساختار درخواست.} تنها فیلد اجباری این مسیر \code{name} است.

\textbf{ساختار پاسخ موفق.} یک شیء از نوع \lr{API Client} با فیلدهای
\code{id}، \code{name}، \code{status}، \code{created\_at} و
\code{updated\_at}.

\begin{latin}
\begin{lstlisting}[style=jsonStyle]
POST /api/v1/api-clients
Authorization: Bearer <SANCTUM_TOKEN>

{
  "name": "Primary Client"
}
\end{lstlisting}
\end{latin}

\listingcaption{نمونه درخواست \code{Create API Client}}

\begin{latin}
\begin{lstlisting}[style=jsonStyle]
{
  "id": 1,
  "name": "Primary Client",
  "status": "active",
  "created_at": "2026-01-29T08:00:00.000000Z",
  "updated_at": "2026-01-29T08:00:00.000000Z"
}
\end{lstlisting}
\end{latin}

\listingcaption{نمونه پاسخ \code{Create API Client}}

\subsection{\lr{GET /api/v1/api-clients/\{apiClient\}/keys}}
\textbf{کارکرد.} دریافت فهرست کلیدهای ثبت‌شده برای یک مشتری \lr{API}.

\textbf{احراز هویت.} توکن \lr{Sanctum}.

\textbf{ساختار درخواست.} شناسه‌ی مشتری در مسیر قرار می‌گیرد و بدنه‌ای
ارسال نمی‌شود.

\textbf{ساختار پاسخ موفق.} پاسخ این مسیر آرایه‌ای از فراداده‌ی کلیدها است.
هر عضو آرایه معمولاً فیلدهای \code{id}، \code{key\_prefix}،
\code{scopes}، \code{rate\_limit\_per\_min}، \code{allowed\_ips}،
\code{expires\_at}، \code{revoked\_at} و \code{created\_at} را دارد.
مقدار واقعی کلید در این مسیر هرگز برگردانده نمی‌شود.

\subsection{\lr{POST /api/v1/api-clients/\{apiClient\}/keys}}
\textbf{کارکرد.} ایجاد کلید \lr{API} جدید برای یک مشتری.

\textbf{احراز هویت.} توکن \lr{Sanctum}.

\textbf{ساختار درخواست.} در این مسیر می‌توان آرایه‌ی اختیاری
\code{scopes} را برای حوزه‌های دسترسی، فیلد
\code{rate\_limit\_per\_min} را برای تعیین سقف درخواست در دقیقه،
\code{allowed\_ips} را برای فهرست آدرس‌های مجاز و \code{expires\_at} را
برای تاریخ انقضای کلید ارسال کرد.

\textbf{ساختار پاسخ موفق.} شیئی شامل \code{api\_key}، \code{key\_prefix}،
\code{expires\_at} و \code{resource}. فیلد \code{api\_key} فقط یک‌بار
بازگردانده می‌شود.

\begin{latin}
\begin{lstlisting}[style=jsonStyle]
POST /api/v1/api-clients/1/keys
Authorization: Bearer <SANCTUM_TOKEN>

{
  "scopes": ["ai:chat", "ai:embeddings"],
  "rate_limit_per_min": 60,
  "allowed_ips": ["127.0.0.1"]
}
\end{lstlisting}
\end{latin}

\listingcaption{نمونه درخواست \code{Create API Key}}

\begin{latin}
\begin{lstlisting}[style=jsonStyle]
{
  "api_key": "gk_abcd1234secretvalue",
  "key_prefix": "abcd1234",
  "expires_at": null,
  "resource": {
    "id": 7,
    "key_prefix": "abcd1234",
    "scopes": ["ai:chat", "ai:embeddings"],
    "rate_limit_per_min": 60,
    "allowed_ips": ["127.0.0.1"],
    "expires_at": null,
    "revoked_at": null,
    "created_at": "2026-01-29T08:00:00.000000Z"
  }
}
\end{lstlisting}
\end{latin}

\listingcaption{نمونه پاسخ \code{Create API Key}}

\subsection{\lr{POST /api/v1/api-keys/\{apiKey\}/revoke}}
\textbf{کارکرد.} لغو کلید \lr{API}.

\textbf{احراز هویت.} توکن \lr{Sanctum}.

\textbf{ساختار درخواست.} این مسیر بدنه ندارد و شناسه‌ی کلید در مسیر ارسال
می‌شود.

\textbf{ساختار پاسخ موفق.} شیئی با فیلد \code{status} و مقدار \code{revoked}.

\subsection{\lr{POST /api/v1/api-keys/\{apiKey\}/rotate}}
\textbf{کارکرد.} چرخش کلید \lr{API} و صدور کلید جدید با ابطال کلید پیشین.

\textbf{احراز هویت.} توکن \lr{Sanctum}.

\textbf{ساختار درخواست.} این مسیر بدنه ندارد.

\textbf{ساختار پاسخ موفق.} همانند مسیر ایجاد کلید، شامل \code{api\_key}،
\code{key\_prefix}، \code{expires\_at} و \code{resource}.

\section{نقاط پایانی هوش مصنوعی}

\subsection{\lr{POST /api/v1/ai/responses}}
\textbf{کارکرد.} ارسال درخواست در قالب \lr{Responses API} سازگار با
\lr{OpenAI}.

\textbf{احراز هویت.} کلید پروژه از طریق \code{Authorization} یا
\code{X-API-Key}.

\textbf{ساختار درخواست.} فیلدهای اجباری \code{provider} و \code{model} و
حداقل یک ورودی در \code{input}. سایر فیلدهای سازگار با \lr{Responses API}
به سمت ارائه‌دهنده عبور داده می‌شوند.

\textbf{ساختار پاسخ موفق.} شیئی از نوع \lr{response} شامل فیلدهایی مانند
\code{id}، \code{object}، \code{model} و \code{output}. در صورت وجود اطلاعات
مصرف، فیلد \code{usage} نیز بازگردانده می‌شود.

\begin{latin}
\begin{lstlisting}[style=jsonStyle]
POST /api/v1/ai/responses
Authorization: Bearer <API_KEY>

{
  "provider": "openai",
  "model": "gpt-4o-mini",
  "input": "Hello"
}
\end{lstlisting}
\end{latin}

\listingcaption{نمونه درخواست \code{Responses API}}

\begin{latin}
\begin{lstlisting}[style=jsonStyle]
{
  "id": "resp_test",
  "object": "response",
  "model": "gpt-4o-mini",
  "output": [
    {
      "type": "message",
      "role": "assistant",
      "content": [
        {
          "type": "output_text",
          "text": "Hello from responses"
        }
      ]
    }
  ]
}
\end{lstlisting}
\end{latin}

\listingcaption{نمونه پاسخ \code{Responses API}}

\subsection{\lr{POST /api/v1/ai/chat/completions}}
\textbf{کارکرد.} ارسال درخواست تکمیل گفتگو.

\textbf{احراز هویت.} کلید پروژه.

\textbf{ساختار درخواست.} فیلدهای اجباری \code{provider}، \code{model} و
\code{messages}. آرایه‌ی \code{messages} باید حداقل یک پیام در قالب
\code{role}/\code{content} داشته باشد.

\textbf{ساختار پاسخ موفق.} شیئی از نوع \code{chat.completion} شامل
\code{id}، \code{object}، \code{created}، \code{model}، \code{choices} و
\code{usage}. این پاسخ هم برای \lr{OpenAI}-مانندها و هم برای \lr{Gemini}
پس از نرمال‌سازی به همین شکل برمی‌گردد.

\begin{latin}
\begin{lstlisting}[style=jsonStyle]
POST /api/v1/ai/chat/completions
Authorization: Bearer <API_KEY>

{
  "provider": "groq",
  "model": "llama-3.3-70b-versatile",
  "messages": [
    {"role": "user", "content": "Hello!"}
  ]
}
\end{lstlisting}
\end{latin}

\listingcaption{نمونه درخواست \code{Chat Completions}}

\begin{latin}
\begin{lstlisting}[style=jsonStyle]
{
  "id": "chatcmpl_groq",
  "object": "chat.completion",
  "created": 1738137600,
  "model": "llama-3.3-70b-versatile",
  "choices": [
    {
      "index": 0,
      "message": {
        "role": "assistant",
        "content": "Hello from Groq"
      },
      "finish_reason": "stop"
    }
  ],
  "usage": {
    "prompt_tokens": 5,
    "completion_tokens": 7,
    "total_tokens": 12
  }
}
\end{lstlisting}
\end{latin}

\listingcaption{نمونه پاسخ \code{Chat Completions}}

\subsection{\lr{POST /api/v1/ai/embeddings}}
\textbf{کارکرد.} ارسال درخواست جاسازی متن.

\textbf{احراز هویت.} کلید پروژه.

\textbf{ساختار درخواست.} فیلدهای اجباری \code{provider}، \code{model} و
\code{input}. مقدار \code{input} می‌تواند رشته یا آرایه‌ای از رشته‌ها باشد.

\textbf{ساختار پاسخ موفق.} شیئی با \code{object=list}، آرایه‌ی \code{data}،
فیلد \code{model} و در صورت وجود، \code{usage}.

\begin{latin}
\begin{lstlisting}[style=jsonStyle]
POST /api/v1/ai/embeddings
X-API-Key: <API_KEY>

{
  "provider": "openai",
  "model": "text-embedding-3-small",
  "input": "hello"
}
\end{lstlisting}
\end{latin}

\listingcaption{نمونه درخواست \code{Embeddings}}

\begin{latin}
\begin{lstlisting}[style=jsonStyle]
{
  "object": "list",
  "data": [
    {
      "object": "embedding",
      "index": 0,
      "embedding": [0.1, 0.2]
    }
  ],
  "model": "text-embedding-3-small",
  "usage": {
    "prompt_tokens": 4,
    "total_tokens": 4
  }
}
\end{lstlisting}
\end{latin}

\listingcaption{نمونه پاسخ \code{Embeddings}}

\subsection{\lr{POST /api/v1/ai/images/generations}}
\textbf{کارکرد.} ارسال درخواست تولید تصویر.

\textbf{احراز هویت.} کلید پروژه.

\textbf{ساختار درخواست.} فیلدهای اجباری \code{provider}، \code{model} و
\code{prompt}. فیلدهای سازگار دیگری مانند \code{n} یا \code{size} نیز
می‌توانند همراه درخواست عبور داده شوند.

\textbf{ساختار پاسخ موفق.} شیئی شامل زمان ایجاد در \code{created} و آرایه‌ی
\code{data} که هر عضو آن معمولاً حاوی \code{url} یا داده‌ی تصویر است.

\begin{latin}
\begin{lstlisting}[style=jsonStyle]
POST /api/v1/ai/images/generations
Authorization: Bearer <API_KEY>

{
  "provider": "openai",
  "model": "gpt-image-1",
  "prompt": "A test image"
}
\end{lstlisting}
\end{latin}

\listingcaption{نمونه درخواست \code{Image Generation}}

\begin{latin}
\begin{lstlisting}[style=jsonStyle]
{
  "created": 1738137600,
  "data": [
    {
      "url": "https://example.com/image.png"
    }
  ]
}
\end{lstlisting}
\end{latin}

\listingcaption{نمونه پاسخ \code{Image Generation}}

\subsection{\lr{POST /api/v1/ai/audio/transcriptions}}
\textbf{کارکرد.} ارسال درخواست تبدیل گفتار به متن.

\textbf{احراز هویت.} کلید پروژه.

\textbf{ساختار درخواست.} فیلدهای اجباری \code{provider} و \code{model} به‌همراه
فایل بارگذاری‌شده در \code{file}. این مسیر باید به‌صورت
\lr{multipart/form-data} فراخوانی شود. بر اساس منطق درگاه، نبود فیلد
\code{file} مستقیماً به خطای \lr{invalid\_request\_error} منجر می‌شود.

\textbf{ساختار پاسخ موفق.} در رایج‌ترین حالت شیئی شامل فیلد \code{text} که
متن رونویسی‌شده را نگه می‌دارد.

\begin{latin}
\begin{lstlisting}[style=jsonStyle]
{
  "text": "hello world"
}
\end{lstlisting}
\end{latin}

\listingcaption{نمونه پاسخ \code{Audio Transcription}}

\subsection{\lr{POST /api/v1/ai/audio/speech}}
\textbf{کارکرد.} ارسال درخواست تبدیل متن به گفتار.

\textbf{احراز هویت.} کلید پروژه.

\textbf{ساختار درخواست.} فیلدهای اجباری \code{provider}، \code{model}،
\code{input} و \code{voice}.

\textbf{ساختار پاسخ موفق.} این مسیر معمولاً \lr{JSON} برنمی‌گرداند، بلکه
پاسخ دودویی صوت به‌همراه \code{Content-Type} مناسب، مانند
\code{audio/mpeg}، تحویل می‌دهد.

\begin{latin}
\begin{lstlisting}[style=jsonStyle]
POST /api/v1/ai/audio/speech
Authorization: Bearer <API_KEY>

{
  "provider": "openai",
  "model": "gpt-4o-mini-tts",
  "input": "Hello there",
  "voice": "alloy"
}
\end{lstlisting}
\end{latin}

\listingcaption{نمونه درخواست \code{Audio Speech}}

\begin{latin}
\begin{lstlisting}[style=jsonStyle]
HTTP/1.1 200 OK
Content-Type: audio/mpeg

<binary audio bytes>
\end{lstlisting}
\end{latin}

\listingcaption{نمونه پاسخ \code{Audio Speech}}

\section{نقاط پایانی مالی}

\subsection{\lr{GET /api/v1/wallet}}
\textbf{کارکرد.} مشاهده‌ی موجودی کیف پول.

\textbf{احراز هویت.} توکن \lr{Sanctum}.

\textbf{ساختار درخواست.} این مسیر بدنه ندارد.

\textbf{ساختار پاسخ موفق.} شیئی با فیلدهای \code{id}، \code{balance}،
\code{currency} و \code{updated\_at}.

\subsection{\lr{POST /api/v1/wallet/topup}}
\textbf{کارکرد.} افزایش موجودی کیف پول در نسخه‌ی فعلی از مسیر داخلی سامانه.

\textbf{احراز هویت.} توکن \lr{Sanctum}.

\textbf{ساختار درخواست.} فیلد اجباری \code{amount} و فیلد اختیاری
\code{reason}.

\textbf{ساختار پاسخ موفق.} شیئی شامل دو بخش \code{wallet} و
\code{transaction}. بخش \code{wallet} موجودی جدید را نشان می‌دهد و بخش
\code{transaction} شامل \code{id}، \code{type}، \code{amount}، \code{reason}،
\code{ref\_type}، \code{ref\_id}، \code{meta} و \code{created\_at} است.

\begin{latin}
\begin{lstlisting}[style=jsonStyle]
POST /api/v1/wallet/topup
Authorization: Bearer <SANCTUM_TOKEN>

{
  "amount": 25.50,
  "reason": "test_topup"
}
\end{lstlisting}
\end{latin}

\listingcaption{نمونه درخواست \code{Wallet Topup}}

\begin{latin}
\begin{lstlisting}[style=jsonStyle]
{
  "wallet": {
    "id": 1,
    "balance": 25.5,
    "currency": "USD",
    "updated_at": "2026-01-29T08:00:00.000000Z"
  },
  "transaction": {
    "id": 1,
    "type": "credit",
    "amount": 25.5,
    "reason": "test_topup",
    "ref_type": null,
    "ref_id": null,
    "meta": null,
    "created_at": "2026-01-29T08:00:00.000000Z"
  }
}
\end{lstlisting}
\end{latin}

\listingcaption{نمونه پاسخ \code{Wallet Topup}}

\subsection{\lr{GET /api/v1/wallet/transactions}}
\textbf{کارکرد.} دریافت دفترکل تراکنش‌های کیف پول.

\textbf{احراز هویت.} توکن \lr{Sanctum}.

\textbf{ساختار درخواست.} پارامتر اختیاری \code{limit} در \lr{query string}
پذیرفته می‌شود.

\textbf{ساختار پاسخ موفق.} آرایه‌ای از اشیای تراکنش با همان ساختار
\code{transaction} در مسیر افزایش موجودی.

\subsection{\lr{GET /api/v1/plans}}
\textbf{کارکرد.} مشاهده‌ی لیست طرح‌های اشتراک.

\textbf{احراز هویت.} توکن \lr{Sanctum}.

\textbf{ساختار درخواست.} این مسیر بدنه ندارد.

\textbf{ساختار پاسخ موفق.} آرایه‌ای از اشیا با فیلدهای \code{id}،
\code{name}، \code{price}، \code{currency}، \code{period}،
\code{included\_credits}، \code{rate\_limits}، \code{features}،
\code{status}، \code{created\_at} و \code{updated\_at}.

\subsection{\lr{POST /api/v1/plans}}
\textbf{کارکرد.} ایجاد طرح اشتراک جدید.

\textbf{احراز هویت.} توکن \lr{Sanctum}.

\textbf{ساختار درخواست.} فیلدهای \code{name}، \code{price} و
\code{period} اجباری هستند. در کنار آن‌ها، فیلدهای \code{currency}،
\code{included\_credits}، \code{rate\_limits}، \code{features} و
\code{status} نیز به‌صورت اختیاری پذیرفته می‌شوند.

\textbf{ساختار پاسخ موفق.} شیء طرح اشتراک با همان فیلدهای مسیر
\code{GET /api/v1/plans}.

\subsection{\lr{POST /api/v1/subscriptions}}
\textbf{کارکرد.} خرید اشتراک.

\textbf{احراز هویت.} توکن \lr{Sanctum}.

\textbf{ساختار درخواست.} تنها فیلد اجباری \code{plan\_id} است.

\textbf{ساختار پاسخ موفق.} شیئی شامل \code{id}، \code{status}،
\code{starts\_at}، \code{ends\_at}، \code{renewal\_at}،
\code{canceled\_at} و شیء تو در توی \code{plan}.

\begin{latin}
\begin{lstlisting}[style=jsonStyle]
POST /api/v1/subscriptions
Authorization: Bearer <SANCTUM_TOKEN>

{
  "plan_id": 1
}
\end{lstlisting}
\end{latin}

\listingcaption{نمونه درخواست \code{Create Subscription}}

\begin{latin}
\begin{lstlisting}[style=jsonStyle]
{
  "id": 25,
  "status": "active",
  "starts_at": "2026-01-29T08:00:00.000000Z",
  "ends_at": "2026-02-29T08:00:00.000000Z",
  "renewal_at": "2026-02-29T08:00:00.000000Z",
  "canceled_at": null,
  "plan": {
    "id": 1,
    "name": "Starter",
    "price": 10,
    "currency": "USD",
    "period": "monthly",
    "included_credits": 100,
    "rate_limits": null,
    "features": null,
    "status": "active",
    "created_at": "2026-01-29T08:00:00.000000Z",
    "updated_at": "2026-01-29T08:00:00.000000Z"
  }
}
\end{lstlisting}
\end{latin}

\listingcaption{نمونه پاسخ \code{Create Subscription}}

\subsection{\lr{GET /api/v1/subscriptions/current}}
\textbf{کارکرد.} دریافت اشتراک فعال کاربر.

\textbf{احراز هویت.} توکن \lr{Sanctum}.

\textbf{ساختار درخواست.} این مسیر بدنه ندارد.

\textbf{ساختار پاسخ موفق.} همان ساختار مسیر \code{POST /api/v1/subscriptions}.
در صورت نبود اشتراک فعال، خطای \code{404} با پیام
\code{No active subscription.} برگردانده می‌شود.

\subsection{\lr{POST /api/v1/subscriptions/cancel}}
\textbf{کارکرد.} لغو اشتراک فعال.

\textbf{احراز هویت.} توکن \lr{Sanctum}.

\textbf{ساختار درخواست.} این مسیر بدنه ندارد.

\textbf{ساختار پاسخ موفق.} همان ساختار منبع اشتراک با \code{status} برابر
\code{canceled}. در صورت نبود اشتراک فعال، خطای \code{404} بازگردانده
می‌شود.

\subsection{\lr{GET /api/v1/invoices}}
\textbf{کارکرد.} مشاهده‌ی لیست فاکتورها.

\textbf{احراز هویت.} توکن \lr{Sanctum}.

\textbf{ساختار درخواست.} پارامتر اختیاری \code{limit} پذیرفته می‌شود.

\textbf{ساختار پاسخ موفق.} آرایه‌ای از خلاصه‌ی فاکتورها با فیلدهای
\code{id}، \code{number}، \code{status}، \code{currency}، \code{subtotal}،
\code{tax}، \code{total}، \code{issued\_at}، \code{paid\_at}،
\code{pdf\_url}، \code{created\_at} و \code{updated\_at}. در مسیر فهرست،
اقلام ریزفاکتور معمولاً بارگذاری نمی‌شوند.

\subsection{\lr{GET /api/v1/invoices/\{invoice\}}}
\textbf{کارکرد.} مشاهده‌ی جزئیات یک فاکتور مشخص.

\textbf{احراز هویت.} توکن \lr{Sanctum}.

\textbf{ساختار درخواست.} شناسه‌ی فاکتور در مسیر قرار می‌گیرد.

\textbf{ساختار پاسخ موفق.} شیء فاکتور با همان فیلدهای مسیر فهرست، به‌علاوه‌ی
آرایه‌ی \code{items}. هر عضو \code{items} شامل \code{id}، \code{type}،
\code{description}، \code{quantity}، \code{unit\_price}،
\code{line\_total}، \code{meta} و \code{created\_at} است.

\begin{latin}
\begin{lstlisting}[style=jsonStyle]
{
  "id": 3,
  "number": "INV-20260129-0003",
  "status": "issued",
  "currency": "USD",
  "subtotal": 10.0,
  "tax": 0.0,
  "total": 10.0,
  "issued_at": "2026-01-29T08:00:00.000000Z",
  "paid_at": null,
  "items": [
    {
      "id": 11,
      "type": "usage",
      "description": "Chat completions usage",
      "quantity": 1500,
      "unit_price": 0.0005,
      "line_total": 0.75,
      "meta": null,
      "created_at": "2026-01-29T08:00:00.000000Z"
    }
  ],
  "pdf_url": "https://example.com/api/v1/invoices/3/pdf?...",
  "created_at": "2026-01-29T08:00:00.000000Z",
  "updated_at": "2026-01-29T08:00:00.000000Z"
}
\end{lstlisting}
\end{latin}

\listingcaption{نمونه پاسخ \code{Invoice Detail}}

\subsection{\lr{GET /api/v1/invoices/\{invoice\}/pdf}}
\textbf{کارکرد.} دریافت نسخه‌ی \lr{PDF} فاکتور از طریق پیوند امضاشده.

\textbf{احراز هویت.} به‌طور معمول از طریق امضای پیوند انجام می‌شود، نه
توکن \lr{Sanctum}.

\textbf{ساختار درخواست.} این مسیر بدنه ندارد و با \lr{signed URL} فراخوانی
می‌شود.

\textbf{ساختار پاسخ موفق.} اگر فایل آماده باشد، پاسخ دودویی \lr{PDF}
برمی‌گردد. اگر هنوز فایل ساخته نشده باشد، مسیر با وضعیت \code{202} شیء
\code{\{"status":"processing"\}} را بازمی‌گرداند و ساخت \lr{PDF} را در
صف اجرا می‌کند.

\section{نقاط پایانی گزارش‌ها}

\subsection{\lr{GET /api/v1/reports/usage}}
\textbf{کارکرد.} گزارش مصرف با امکان گروه‌بندی و فیلتر بازه‌ی زمانی.

\textbf{احراز هویت.} توکن \lr{Sanctum}.

\textbf{ساختار درخواست.} پارامترهای اختیاری \code{from}، \code{to} و
\code{group\_by}. مقدار \code{group\_by} فقط می‌تواند یکی از
\code{day}، \code{key} یا \code{provider} باشد.

\textbf{ساختار پاسخ موفق.} آرایه‌ای از ردیف‌ها با فیلدهای \code{group\_by}،
\code{group\_value}، \code{metric}، \code{quantity} و \code{total\_cost}.

\begin{latin}
\begin{lstlisting}[style=jsonStyle]
[
  {
    "group_by": "day",
    "group_value": "2026-01-29",
    "metric": "tokens_in",
    "quantity": 100,
    "total_cost": 0.01
  }
]
\end{lstlisting}
\end{latin}

\listingcaption{نمونه پاسخ \code{Usage Report}}

\subsection{\lr{GET /api/v1/reports/wallet-ledger}}
\textbf{کارکرد.} دفترکل تراکنش‌های کیف پول در یک بازه‌ی زمانی.

\textbf{احراز هویت.} توکن \lr{Sanctum}.

\textbf{ساختار درخواست.} پارامترهای اختیاری \code{from} و \code{to}.

\textbf{ساختار پاسخ موفق.} آرایه‌ای از اشیای دفترکل با فیلدهای \code{id}،
\code{type}، \code{amount}، \code{reason}، \code{ref\_type}،
\code{ref\_id}، \code{meta} و \code{created\_at}.

\subsection{\lr{GET /api/v1/reports/invoices}}
\textbf{کارکرد.} گزارش تجمیعی فاکتورها و وضعیت آن‌ها.

\textbf{احراز هویت.} توکن \lr{Sanctum}.

\textbf{ساختار درخواست.} پارامتر اختیاری \code{status} که یکی از مقادیر
\code{draft}، \code{issued}، \code{paid} یا \code{void} را می‌پذیرد.

\textbf{ساختار پاسخ موفق.} آرایه‌ای از خلاصه‌ی فاکتورها با همان ساختار
مسیر \code{GET /api/v1/invoices}.

\section{نقاط پایانی مدیریت ارائه‌دهندگان}

\subsection{\lr{GET /api/v1/providers}}
\textbf{کارکرد.} دریافت فهرست ارائه‌دهندگان ثبت‌شده در سامانه.

\textbf{احراز هویت.} توکن \lr{Sanctum}.

\textbf{ساختار درخواست.} این مسیر بدنه ندارد.

\textbf{ساختار پاسخ موفق.} آرایه‌ای از اشیا با فیلدهای \code{id}،
\code{name}، \code{type}، \code{base\_url}، \code{status}،
\code{priority}، \code{created\_at} و \code{updated\_at}.

\subsection{\lr{POST /api/v1/providers}}
\textbf{کارکرد.} ایجاد ارائه‌دهنده‌ی جدید و ثبت پیکربندی اولیه.

\textbf{احراز هویت.} توکن \lr{Sanctum}.

\textbf{ساختار درخواست.} در این مسیر، فیلدهای \code{name} و
\code{base\_url} اجباری هستند و فیلدهای \code{type}، \code{api\_key}،
\code{timeout}، \code{status} و \code{priority} می‌توانند به‌صورت
اختیاری همراه درخواست ارسال شوند.

\textbf{ساختار پاسخ موفق.} یک شیء ارائه‌دهنده با همان فیلدهای مسیر
\code{GET /api/v1/providers}.

\begin{latin}
\begin{lstlisting}[style=jsonStyle]
POST /api/v1/providers
Authorization: Bearer <SANCTUM_TOKEN>

{
  "name": "openrouter",
  "type": "openai_compatible",
  "base_url": "https://openrouter.ai/api/v1",
  "api_key": "test-key",
  "status": "active"
}
\end{lstlisting}
\end{latin}

\listingcaption{نمونه درخواست \code{Create Provider}}

\begin{latin}
\begin{lstlisting}[style=jsonStyle]
{
  "id": 4,
  "name": "openrouter",
  "type": "openai_compatible",
  "base_url": "https://openrouter.ai/api/v1",
  "status": "active",
  "priority": 0,
  "created_at": "2026-01-29T08:00:00.000000Z",
  "updated_at": "2026-01-29T08:00:00.000000Z"
}
\end{lstlisting}
\end{latin}

\listingcaption{نمونه پاسخ \code{Create Provider}}

\subsection{\lr{GET /api/v1/providers/\{provider\}/models}}
\textbf{کارکرد.} دریافت مدل‌های ثبت‌شده برای یک ارائه‌دهنده.

\textbf{احراز هویت.} توکن \lr{Sanctum}.

\textbf{ساختار درخواست.} شناسه‌ی عددی ارائه‌دهنده در مسیر قرار می‌گیرد.

\textbf{ساختار پاسخ موفق.} آرایه‌ای از اشیا با فیلدهای \code{id}،
\code{provider\_id}، \code{model\_key}، \code{capabilities}،
\code{pricing\_config}، \code{status}، \code{created\_at} و
\code{updated\_at}.

\subsection{\lr{POST /api/v1/providers/\{provider\}/models}}
\textbf{کارکرد.} ایجاد یا ثبت مدل جدید برای یک ارائه‌دهنده.

\textbf{احراز هویت.} توکن \lr{Sanctum}.

\textbf{ساختار درخواست.} در این مسیر \code{model\_key} اجباری است.
فیلدهای \code{capabilities}، \code{pricing\_config} و \code{status}
به‌صورت اختیاری دریافت می‌شوند.

\textbf{ساختار پاسخ موفق.} یک شیء مدل با همان فیلدهای مسیر فهرست مدل‌ها.

\begin{latin}
\begin{lstlisting}[style=jsonStyle]
POST /api/v1/providers/4/models
Authorization: Bearer <SANCTUM_TOKEN>

{
  "model_key": "gpt-4o-mini",
  "capabilities": {
    "chat": true
  },
  "pricing_config": {
    "input_token": 0.0005,
    "output_token": 0.001
  }
}
\end{lstlisting}
\end{latin}

\listingcaption{نمونه درخواست \code{Create Provider Model}}

\begin{latin}
\begin{lstlisting}[style=jsonStyle]
{
  "id": 9,
  "provider_id": 4,
  "model_key": "gpt-4o-mini",
  "capabilities": {
    "chat": true
  },
  "pricing_config": {
    "input_token": 0.0005,
    "output_token": 0.001
  },
  "status": "active",
  "created_at": "2026-01-29T08:00:00.000000Z",
  "updated_at": "2026-01-29T08:00:00.000000Z"
}
\end{lstlisting}
\end{latin}

\listingcaption{نمونه پاسخ \code{Create Provider Model}}

\section{نقاط پایانی ممیزی}

\subsection{\lr{GET /api/v1/audit-logs}}
\textbf{کارکرد.} دریافت لاگ عملیات حساس برای کاربران دارای نقش مدیر.

\textbf{احراز هویت.} توکن \lr{Sanctum} به‌همراه عبور از میان‌افزار
\code{role:admin}.

\textbf{ساختار درخواست.} پارامترهای اختیاری \code{actor\_user\_id}،
\code{action}، \code{target\_type}، \code{target\_id} و \code{limit}.

\textbf{ساختار پاسخ موفق.} آرایه‌ای از اشیای ممیزی با فیلدهای \code{id}،
\code{actor\_user\_id}، \code{action}، \code{target\_type}،
\code{target\_id}، \code{meta} و \code{created\_at}.





